\documentclass[11pt]{article}
\usepackage[a4paper, top=18mm, bottom=18mm, left=20mm, right=20mm]{geometry}
\usepackage[T2A]{fontenc}
\usepackage[utf8]{inputenc}
\usepackage[ukrainian]{babel}
\usepackage{graphicx}
\usepackage[export]{adjustbox}
\usepackage{amsmath}
\usepackage{amssymb}
\setlength{\parindent}{0pt}
\setlength{\parskip}{4pt}
\pagestyle{empty}
\begin{document}
%begin
{\large\textbf{Контрольна робота}}

\textbf{Варіант \No\,1}\hfill \textit{ПІБ:}\,\underline{\hspace{60mm}}
\vspace{4mm}

%beginitem
1. Що буде виведено за виконання фрагменту коду?

%code
cout << (12 / 12 / 12 / 12);
%code

%enditem
%beginitem
2. %goodpath:Scripts/2019/Mod2019_1/Lex/01-good.txt
Відмітити все, що є однією правильною лексемою:
( 1 ) vsop

( 2 ) throw

( 3 ) double

( 4 ) f"123"

%enditem


%end
\newpage

\newpage

%begin
{\large\textbf{Контрольна робота}}

\textbf{Варіант \No\,2}\hfill \textit{ПІБ:}\,\underline{\hspace{60mm}}
\vspace{4mm}

%beginitem
1. %goodpath:Scripts/2018/Mod2018_1/03-good.txt
Відмітити все, що є коректним оператором С++:
( 1 ) x+=y;

( 2 ) if ('x'><'y') a=6;

( 3 ) 'a'=3;

( 4 ) if ('x'<'y') a=6;

%enditem
%beginitem
2. Що буде виведено за виконання фрагменту коду?

%code
cout << (!8.0<=9.0 or !8>= 9 or false>4);
%code

%enditem


%end
\newpage

\newpage

%begin
{\large\textbf{Контрольна робота}}

\textbf{Варіант \No\,3}\hfill \textit{ПІБ:}\,\underline{\hspace{60mm}}
\vspace{4mm}

%beginitem
1. 
try:
    res = 4**-0.5*-1
    if isinstance(res, bool):
        print(f'bool;{res}')
    elif isinstance(res, int):
        print(f'int;{res}')
    elif isinstance(res, float):
        print(f'float;{res:.2f}')
    else:
        print('???')
except:
    print('Z')

%enditem
%beginitem
2. %goodpath:Scripts/2019/Mod2019_1/Lex/03-good.txt
Відмітити все, що є коректним оператором С++:
( 1 ) "\n" is not '\n'

( 2 ) x in not y

( 3 ) (2+2j)%2

( 4 ) print(3%50)

%enditem


%end
\newpage

\newpage

%begin
{\large\textbf{Контрольна робота}}

\textbf{Варіант \No\,4}\hfill \textit{ПІБ:}\,\underline{\hspace{60mm}}
\vspace{4mm}

%beginitem
1. %goodpath:Scripts/2019/Mod2019_1/Lex/03-good.txt
Відмітити все, що є коректним оператором С++:
( 1 ) "\n" is not '\n'

( 2 ) x in not y

( 3 ) (2+2j)%2

( 4 ) print(3%50)

%enditem
%beginitem
2. 
a = 6
b = 7
c = 8

def g(b):
global c    a -= 5
    b = 3
    c = 1
    return a + b + c

a = 4
b = 9
c = 0

try:
    print(g(a), a, b, c, sep=';')
except:
    print('Z', a, b, c, sep=';')

%enditem


%end
\newpage

\newpage

%begin
{\large\textbf{Контрольна робота}}

\textbf{Варіант \No\,5}\hfill \textit{ПІБ:}\,\underline{\hspace{60mm}}
\vspace{4mm}

%beginitem
1. 

print('R', end=' ')
j=6
while j<=12:
    j+=2
print(j)


%enditem
%beginitem
2. %goodpath:Scripts/2018/Mod2018_1/01-good.txt
Відмітити все, що є однією правильною лексемою:
( 1 ) "1.1.1"

( 2 ) (1)

( 3 ) 2int

( 4 ) "1,3,2"

%enditem


%end
\newpage

\newpage

%begin
{\large\textbf{Контрольна робота}}

\textbf{Варіант \No\,6}\hfill \textit{ПІБ:}\,\underline{\hspace{60mm}}
\vspace{4mm}

%beginitem
1. Що буде виведено за виконання фрагменту коду?

%code
print('R', end=' ')
j=12
while j>=6:
    j+=-2
print(j)
%code


%enditem
%beginitem
2. %goodpath:Scripts/2018/Mod2018_1/03-good.txt
Відмітити все, що є коректним оператором С++:
( 1 ) x+=y;

( 2 ) if ('x'><'y') a=6;

( 3 ) 'a'=3;

( 4 ) if ('x'<'y') a=6;

%enditem


%end
\newpage

\newpage

%begin
{\large\textbf{Контрольна робота}}

\textbf{Варіант \No\,7}\hfill \textit{ПІБ:}\,\underline{\hspace{60mm}}
\vspace{4mm}

%beginitem
1. %goodpath:Scripts/2018/Mod2018_1/03-good.txt
Відмітити все, що є коректним оператором С++:
( 1 ) x+=y;

( 2 ) if ('x'><'y') a=6;

( 3 ) 'a'=3;

( 4 ) if ('x'<'y') a=6;

%enditem
%beginitem
2. Що буде виведено за виконання фрагменту коду?

%code
a,b,c=6,7,8
def g(b):
    a=5
    b=3
    c=1
    return a+b+c

a,b,c=4,9,0
print(g(a),a,b,c)
%code

%enditem


%end
\newpage

\newpage

%begin
{\large\textbf{Контрольна робота}}

\textbf{Варіант \No\,8}\hfill \textit{ПІБ:}\,\underline{\hspace{60mm}}
\vspace{4mm}

%beginitem
1. %goodpath:Scripts/2018/Mod2018_1/01-good.txt
Відмітити все, що є однією правильною лексемою:
( 1 ) "1.1.1"

( 2 ) (1)

( 3 ) 2int

( 4 ) "1,3,2"

%enditem
%beginitem
2. Що буде виведено за виконання фрагменту коду?

%code
print(7//7/7*7)
%code

%enditem


%end
\newpage

\newpage

%begin
{\large\textbf{Контрольна робота}}

\textbf{Варіант \No\,9}\hfill \textit{ПІБ:}\,\underline{\hspace{60mm}}
\vspace{4mm}

%beginitem
1. Що буде виведено за виконання фрагменту коду?

%code
print('R', end=' ')
g=['0','1','2','3','4']
g[-3:]=[2,3]
print(g)
%code


%enditem
%beginitem
2. %goodpath:Scripts/2018/Mod2018_1/03-good.txt
Відмітити все, що є коректним оператором С++:
( 1 ) x+=y;

( 2 ) if ('x'><'y') a=6;

( 3 ) 'a'=3;

( 4 ) if ('x'<'y') a=6;

%enditem


%end
\newpage

\newpage

%begin
{\large\textbf{Контрольна робота}}

\textbf{Варіант \No\,10}\hfill \textit{ПІБ:}\,\underline{\hspace{60mm}}
\vspace{4mm}

%beginitem
1. %goodpath:Scripts/2019/Mod2019_1/Lex/03-good.txt
Відмітити все, що є коректним оператором С++:
( 1 ) "\n" is not '\n'

( 2 ) x in not y

( 3 ) (2+2j)%2

( 4 ) print(3%50)

%enditem
%beginitem
2. Що буде виведено за виконання фрагменту коду?

%code
print(7//7/7*7)
%code

%enditem


%end
\newpage
\end{document}


